\section{Related Work}
\label{sec:related}

\paragraph{Graph retrieval-augmented generation.}
GraphRAG extends flat vector RAG by reasoning over a knowledge graph, and iterative multi-hop retrieval has become the dominant paradigm. Think-on-Graph (ToG)~\cite{sun2023thinkongraph} introduced LLM-guided beam search over a KG, expanding a frontier of triples and pruning by LLM relevance at each hop; ToG~2.0~\cite{ma2024thinkongraph} added knowledge-guided retrieval, and Search-on-Graph~\cite{sun2025searchongraph} refined iterative informed navigation. GNN-RAG~\cite{mavromatis2025gnnrag} decouples graph reasoning from generation by training a GNN to score candidate paths, while Paths-over-Graph~\cite{tan2025pathsovergraph} and Fast-ToG~\cite{liang2025fast} pursue wider or faster reasoning. These methods share a fixed retrieval budget applied uniformly to every query, which is the inefficiency \method{} targets. Surveys~\cite{peng2024graph,zhang2025survey,huang2026survey} confirm the field's rapid growth and the emerging emphasis on cost.

\paragraph{Adaptive and cost-efficient GraphRAG.}
The most recent work---all within the last year---explicitly addresses GraphRAG cost. \emph{Use-Graph-When-It-Needs}~\cite{dong2026graph} learns when to invoke the graph, and \emph{GraphRAG-Router}~\cite{fan2026graphragrouter} learns a reinforcement-learned routing policy over backends. Core-based Hierarchies~\cite{hossain2026corebased} and EHRAG~\cite{song2026ehrag} reduce indexing cost, while S-Path-RAG~\cite{fu2026spathrag} and the minimal-subgraph retriever of Yuan et al.~\cite{yuan2026retrieving} prune the retrieval path. These approaches either \emph{train} a router/detector (adding a learning stage and a domain-specific policy) or optimize the \emph{index} rather than per-query \emph{retrieval budget}. \method{} differs by being fully training-free and by adapting the retrieval budget \emph{at inference time} using signals already computed by the pipeline. SkewRoute~\cite{wang2025skewroute} is closest in spirit---it routes using score-skewness---but targets LLM-backend selection, not graph-retrieval budgeting.

\paragraph{Training-free adaptation and early stopping.}
The idea of adapting inference compute to input difficulty without training appears in adaptive reasoning-suppression for reasoning LMs~\cite{zheng2025adaptive}, token-efficient prompting~\cite{shah2026crop}, and early-exit transformers. \method{} applies this principle to \emph{retrieval} rather than generation: the controller adapts the number of retrieval hops and the beam width, not the generation length. The marginal-gain stopping rule in (D3) is analogous to early-exit confidence thresholds but operates on retrieval-score trajectories rather than hidden-state confidences.

\paragraph{Knowledge-graph question answering.}
MetaQA~\cite{zhang2017variational} is the canonical multi-hop KGQA benchmark, with earlier neural approaches including EmbedKGQA~\cite{saxena2020improving}, PullNet~\cite{sun2019pullnet}, ReaRev~\cite{mavromatis2022rearev}, and TransferNet~\cite{shi2021transfernet}. The LLM era shifted KGQA from learned graph embeddings to LLM-guided retrieval~\cite{sun2023thinkongraph,li2023fewshot,mavromatis2025gnnrag}. We adopt MetaQA for comparability with this line and report F1/Hit@1 under the same protocol.
